% !TeX root = ./0-PhD-Thesis-JLBG.tex
\label{cap: list of publications}
\hypertarget{listofpublications}{During} my doctoral studies, I contributed to the publication of the following peer-reviewed works:
%
\vspace{0.5cm}

%
%
\begin{mybox2}{\centering 1.- Angular momentum transfer from swift electron to small non-spherical nanoparticles in the dipolar approximation  (2024)}
	\textbf{J. L. Briseño-Gómez}, A. López-Tercero, J. Á. Castellanos-Reyes, A. Reyes-Coronado,
	\href{https://doi.org/10.1016/j.ultramic.2024.114005}{Ultramicroscopy,
		Volume 264,
		114005 
		(2024).}
\end{mybox2}{}
%
In electron microscopy and material science, the interaction between swift electrons and nanostructures plays a crucial role. In this work, we present a comprehensive theoretical investigation into the influence of nanoparticle geometry on angular momentum transfer from a swift electron to small symmetric nanoparticles. Specifically, we focus on spheroidal and polyhedral nanoparticles, analyzing their response within the framework of the dipolar approximation. By exploring various geometric configurations, we showed how subtle changes in shape affect the efficiency and dynamics of angular momentum transfer. This study not only enhances our understanding of the underlying physical mechanisms but also provides valuable insights for potential applications, such as nanomanipulation or fabrication of nanomaterials. 

\vspace{0.5cm}

%Además, en este momento se encuentra en proceso de revisión el trabajo (EJEMPLO)

%\vspace{0.5cm}
%
%
%\begin{mybox2}{\centering 4.- Non-causality and numerical convergence impact on the linear momentum transfer by a swift electron to a metallic nanoparticle}
%	J. Castrejón-Figueroa, \textbf{J. Á. Castellanos-Reyes}, A. Reyes-Coronado, 
%	\href{https://arxiv.org/abs/2106.04032}{arXiv:2106.04032.}
%\end{mybox2}{}
%
%En este trabajo se presenta un estudio teórico sobre los efectos de la no causalidad y la falta de convergencia numérica en los cálculos de la transferencia de momento lineal de haces de electrones a nanopartículas.
%\vspace{0.5cm}

We are also preparing two manuscripts for submission as peer-reviewed publications:
%
\begin{itemize}
	\item The first manuscript, authored by \textbf{J. L. Briseño-Gómez}, A. Reyes-Coronado, and M. Á. Pizaña-Juárez, focuses on a unique academic and educational-focused problem solution. It involves calculating the electromotive force hypothetically generated on the Moon due to its natural motion through the surrounding spatial magnetic field. This study explores theoretical aspects of electromagnetism in a novel context and aims to provide a pedagogical tool for teaching complex physical concepts and numerical methods.
	\item The second manuscript, co-authored by \textbf{J. L. Briseño-Gómez}, J. Castrejón-Figueroa, E. E. Viveros-Armas, and A. Reyes-Coronado, addresses the efficient calculation of linear momentum transfer from swift electrons to spherical nanoparticles. This work spans the entire nanoscale and examines both plasmonic and dielectric materials using a causal dielectric function, providing a comprehensive study on electron-nanoparticle interactions that could have implications for advanced nanotechnology applications.	
	\item The main manuscript, co-authored by \textbf{J. L. Briseño-Gómez} and A. Reyes-Coronado, focuses on the efficient calculation of angular momentum transfer from swift electrons to spherical nanoparticles. The preliminary results presented in this manuscript are drawn from the work outlined in this PhD Candidate Protocol. 
\end{itemize}

